\begin{table}[tbp]
    \centering
\small
    \setlength{\tabcolsep}{6pt}
    \begin{tabular}{@{}lrrrrr@{}}
        \toprule
        \textbf{Probe stage} & \textbf{Candidates [bits]} & \textbf{Controls [bits]} & \textbf{Gap [bits]} & \textbf{Cells} & \textbf{Per cell} \\
        \midrule
        encoder block $0$  & $78.9$ & $81.2$ & $-2.4$  & $21$ & $80$ \\
        encoder block $1$  & $75.7$ & $71.5$ & $+4.2$  & $21$ & $80$ \\
        encoder block $2$  & $78.0$ & $82.4$ & $-4.4$  & $21$ & $80$ \\
        encoder block $3$  & $88.1$ & $94.9$ & $-6.8$  & $21$ & $80$ \\
        \addlinespace[2pt]
        decoder block $0$  & $89.9$ & $77.7$ & $+12.3$ & $21$ & $80$ \\
        decoder block $1$  & $71.1$ & $58.2$ & $+12.9$ & $21$ & $80$ \\
        decoder block $2$  & $84.9$ & $60.0$ & $+24.8$ & $21$ & $80$ \\
        decoder block $3$  & $85.2$ & $65.4$ & $+19.8$ & $21$ & $80$ \\
        \addlinespace[2pt]
        pre-projection     & $85.7$ & $60.7$ & $+25.0$ & $21$ & $80$ \\
        output head        & $75.8$ & $64.3$ & $+11.5$ & $21$ & $80$ \\
        \bottomrule
    \end{tabular}
\caption{The numeric form of \autoref{fig:appFdepth}: mean local codelength in bits per probe stage at the retrained reference geometry, for candidate frequencies and for their controls, with the gap as candidates minus controls. Each entry averages $21$ cells of $80$ examples. The sign of the gap is not stable across the encoder, three of its four blocks charging less for a candidate than for a control, and turns positive at the first decoder block and stays positive to the head. The spread across cells, drawn as bands in the figure, overlaps at every stage, so these means locate where the cost appears and do not establish that it is there.}
    \label{tab:appFprofile}
\end{table}
